\section{Introduction}
\label{sec:intro}
% The topic of Novel-view synthesis (NVS) has seen impressive progress due to the use of neural networks to learn
% representations that are well suited for view synthesis tasks, known as Neural Radiance Field \cite{nerf}, which fulfill success in various tasks\cite{munkberg2022NVDiffrec,wang2021nerf--,zhang2020nerf++,martin2021nerfwild,barron2021mipnerf}.After the emergence of NeRF, researchers immediately applied it to inverse rendering and relighting tasks
% Real-time has been a longstanding pursuit in the field of computer graphics. Since the last century, researchers have been devoted to addressing issues such as real-time shadowing \cite{chin1989near,guennebaud2006real}, global illumination \cite{pan2007precomputed,sloan2002precomputed}, ray tracing \cite{woop2005rpu,nomoto2020realistic}, and multi-light illumination \cite{bitterli2020spatiotemporal,ouyang2021restir}. 
3D Gaussian Splatting \cite{3dgs} has garnered significant attention from the community as a promising approach for various tasks in 3D scene reconstruction. 
The utilization of 3DGS presents the potential for individuals to reconstruct their surrounding environment using contemporary technological devices such as smartphones and computers in minutes. Furthermore, individuals can modify their reconstructed world according to their unique preferences, which makes it particularly attractive for multimedia applications and could potentially spur innovation within the multimedia industry. The foundation for achieving this lies in a real-time and high-quality inverse rendering, relighting, and scene editing method. 

The achievement of real-time inverse rendering and relighting has been a long-standing problem. Methods based on Neural Radiance Fields (NeRF) \cite{jin2023tensoir,munkberg2022NVDiffrec,zhang2021nerfactor} have exhibited noteworthy accomplishments in high-quality material editing, illumination computing, and shadow estimation. However, these techniques struggle with the computational overhead and cannot achieve the desired quality in dynamic environmental lighting conditions. The integration of MLPs within these methods gives rise to inherent obstacles owing to their restricted expressive capacity and substantial computational demands. These obstacles engender a considerable curtailment of the effectiveness and efficiency of inverse rendering. Current works \cite{shi2023gir,liang2023gsir,gao2023relightable} introduce 3DGS to inverse rendering instead of NeRF, achieving high-performance inverse rendering and relighting by employing a set of 3D Gaussian splats to represent a 3D scene. However, they still face challenges in the real-time computation of high-quality indirect lighting and shadows in dynamic environmental lighting conditions, primarily due to the utilization of ray tracing or ambient occlusion techniques. The former struggles with balancing quality and speed because of the immense number of Gaussian splats involved \cite{gao2023relightable}, while the latter have difficulty in computing realistic indirect lighting \cite{liang2023gsir}. 

% Real-time has been a longstanding pursuit in the field of computer graphics. Since the last century, researchers have been devoted to addressing issues such as real-time shadowing, global illumination, ray tracing, and multi-light illumination.

In this paper, we solve the aforementioned challenge by introducing Precomputed Radiance Transfer (PRT) \cite{sloan2002precomputed} to 3DGS. Starting with assigning each 3D Gaussian splat with normal (geometry), visibility, and BRDF attributes, we precompute the expensive transport simulation required by complex transfer functions like shadowing and interreflection for given 3d Gaussian splats. The precomputed transfer functions and incident radiance are encoded as either a dense set of vectors (diffuse cases) or matrices (glossy cases), utilizing a low-order spherical harmonic (SH) basis for each Gaussian splats distribution. This approach allows for an efficient representation of the complex transfer functions while maintaining a high level of accuracy. Leveraging the linearity inherent in light transport, we streamline the light integration process to a straightforward dot product operation between their coefficient vectors for diffuse surfaces, or a compact transfer matrix for glossy surfaces, significantly reducing computational overhead. 

Our approach not only enhances real-time rendering quality in dynamic lighting but also opens up new possibilities for interactive multimedia applications. We developed unique precomputing and ray tracing methods that are specifically tailored to accommodate the unique geometry and rendering pipeline of 3DGS. This ensures that our approach is optimized to provide high-quality results while maintaining computational efficiency. Specifically, during the training stage, we carefully designed different reflection strategies for various surfaces under different conditions, ensuring the quality of interreflection while significantly reducing training time. Our unique ray tracing method allows us to perform only one-bounce ray tracing throughout the whole training or testing stage, resulting in real-time photorealistic rendering results. Extensive experimental analyses have demonstrated that our proposed approach significantly outperforms existing methods on synthetic and real-world datasets across multiple tasks. We summarize our main contributions as follows:
\begin{itemize}
     \vspace{-0.4cm}
\item[$\bullet$]  We applied PRT to 3DGS for the first time, achieving high-quality real-time relighting in complex scenes under dynamic lighting conditions, while also supporting high-quality scene editing.


\item[$\bullet$] For efficiency, we designed distinct precomputing methods for both training and rendering stage. Additionally, we devised unique ray tracing and indirect lighting precomputation methods for 3DGS to accelerate training speed and compute accurate indirect illumination related to environmental lighting.
\item[$\bullet$] 
Through comprehensive experimentation, we have demonstrated that our approach outperforms relevant schemes. The experimental results highlight that our method not only facilitates high-quality real-time relighting but also excels in supporting high-quality scene editing.
\end{itemize}
%2.我们首次使用新一代相机构建了低光驾驶脉冲流数据集,xxx。它是比现有低光脉冲流数据集具有更大的规模更丰富的场景。
%3.实验结果展示了我们的增强框架可以既可以making the reconstructed results more perceptible and 还能改善其他下游任务的性能。


%去掉三 把1拆成两部分讲 网络与解码器
